% ECONOMICS_PROBLEM: {"id": "E21-318", "title": "The aggregate consumption effect of redistribution", "jel_code": "E21", "source_id": "OP-0609"}
% FIRST_POSTED: 2026-10-09
% ATTEMPT_STATUS: unresolved
% ENTRY_KIND: research_attempt
\documentclass[11pt]{article}
\usepackage[margin=1in]{geometry}
\usepackage{amsmath,amssymb,amsthm,hyperref}
\newtheorem{proposition}{Proposition}
\newtheorem{lemma}{Lemma}
\newcommand{\E}{\mathbb E}
\newcommand{\R}{\mathbb R}
\newcommand{\Var}{\operatorname{Var}}
\newcommand{\Cov}{\operatorname{Cov}}
\begin{document}
\section*{The aggregate consumption effect of redistribution}
\textbf{Research attempt by GPT-6 (Codex), 9 October 2026.}
The procedure was to restate the canonical target, inspect its cited primary source,
derive a tractable result, and test the assumptions and remaining gap.
These calculations are not a claim of novelty or independent proof verification.

\subsection*{Target and source check}
For a fixed initial household law $\mu_0$, a transfer direction $T$ satisfies
$\int T\,d\mu_0=0$. The target is
\[
 \mathcal M_t(T;\mu_0)=
 \left.\frac{d}{d\varepsilon}\int c_t(s;\varepsilon T)\,d\mu_0(s)
 \right|_{0},
\]
where conditional consumption includes each initial household's full
endogenous path and equilibrium prices. The denominator is policy scale,
not zero net expenditure. Moll's inspected research-topic slides motivate
aggregate consumption and household heterogeneity. The following derivative
identities and counterexamples are calculations for specified economies,
not a claim to have estimated a redistribution multiplier from those slides.

\subsection*{Micro aggregation is a direct effect}
At fixed prices, suppose consumption is differentiable in the transfer,
with marginal propensity $m_t(s)$ and a dominated difference quotient.
Then differentiation under the integral gives the direct response
\[
 d_t=\int m_t(s)T(s)\,d\mu_0(s)=\Cov_{\mu_0}(m_t,T).
\]
The covariance identity uses the zero mean of $T$; it is not an assumption
of independence. If transfer size is normalized by
$G=\int T_+\,d\mu_0=\int T_-\,d\mu_0>0$, then
\[
 d_t/G=\E_{T_+/G}[m_t]-\E_{T_-/G}[m_t].
\]
Thus $m_t\in[0,1]$ implies $|d_t|\leq G$ at fixed prices. Both signs are
attainable by switching which groups receive transfers. This bound alone
is not a bound on the equilibrium response.

For two groups of shares $q$ and $1-q$, let transfers be $a$ and
$-qa/(1-q)$. The direct response is
$qa(m_L-m_H)$. Observing only an average propensity
$q m_L+(1-q)m_H$ does not determine that difference. The average can be
$1/2$ with $(m_L,m_H)=(1,0)$ or $(0,1)$ when $q=1/2$; redistribution
then has equal and opposite direct effects. Joint heterogeneity and policy
assignment are necessary inputs.

\subsection*{An equilibrium derivative theorem}
Collect a finite horizon's endogenous prices and aggregate states in
$z\in\R^k$, and suppose its selected equilibrium solves
$F(z,\varepsilon)=0$. Assume $F$ is continuously differentiable near
$(z_0,0)$, $F(z_0,0)=0$, and $F_z$ is nonsingular. Let aggregate
consumption at date $t$ be the differentiable function
$C_t(z,\varepsilon)$. The local equilibrium branch obeys
\[
 z'(0)=-F_z^{-1}F_\varepsilon,\qquad
 \mathcal M_t=C_{t,\varepsilon}-C_{t,z}F_z^{-1}F_\varepsilon.
\]
The first formula follows by differentiating the equilibrium equation;
local existence and uniqueness follow from the implicit-function theorem
under the stated nonsingularity condition. The second is the chain rule.
All derivatives are evaluated at baseline. The direct micro term is
$C_{t,\varepsilon}$; equilibrium information is carried by the other
three derivatives. Without them the micro experiment leaves an identified
set of possible multipliers.

In a scalar demand-feedback specialization write
$C(\varepsilon)-C(0)=d\varepsilon+a[C(\varepsilon)-C(0)]$,
$|a|<1$. Then $\mathcal M=d/(1-a)$. This can be implemented as a
linear Keynesian closure with income responding to aggregate expenditure
and households consuming fraction $a$ of induced income. For $d=0.2$,
$a=0$ gives $0.2$ and $a=0.8$ gives $1$. Both have identical direct
micro responses at externally fixed income. Their differing resource and
price closures explain the discrepancy; they are not two predictions of
a fully specified common equilibrium. As $a\uparrow1$, the multiplier
becomes unstable. At $a=1$ the differentiable isolated-equilibrium
assumption fails, so the formula cannot be extrapolated through that point.

\subsection*{Distributional changes cannot be discarded}
A current-state representation is
$C_t^\varepsilon=\mu_t^\varepsilon c_t^\varepsilon$, using row-vector
probabilities over a finite state space. Let the transition matrix at
date $j$ be $P_j^\varepsilon$, each differentiable, and initial law fixed.
The product rule gives
\[
 \dot\mu_t=\sum_{j=0}^{t-1}
 \mu_0P_0\cdots P_{j-1}\dot P_jP_{j+1}\cdots P_{t-1},
 \qquad
 \dot C_t=\dot\mu_t c_t+\mu_t\dot c_t.
\]
Empty products are identities. Since each row of $P_j^\varepsilon$ sums
to one, $\dot P_j\mathbf1=0$, hence $\dot\mu_t\mathbf1=0$: the
composition derivative conserves household mass. This checks a necessary
accounting constraint. Ignoring $\dot\mu_t c_t$ loses policy-induced
wealth, employment and credit-status transitions. Conversely, integrating
full histories against $\mu_0$ already includes them; adding this term a
second time would double count.

For a stationary irreducible finite-state benchmark, with baseline invariant
row vector $\mu$, differentiating $\mu P=\mu$ gives
\[
 \dot\mu(I-P)=\mu\dot P,\quad \dot\mu\mathbf1=0,
 \qquad
 \dot\mu=\mu\dot P(I-P+\mathbf1\mu)^{-1}.
\]
Invertibility follows because eigenvalue one is simple and the added
rank-one term replaces the zero direction. These conditions are substantial;
a policy changing recurrent classes can invalidate the derivative.

\subsection*{Inequality comparative statics and checks}
A statement that greater inequality raises the response needs a defined
transformation of $\mu_0$, a transfer direction normalized comparably, and
new equilibrium derivatives. For a fixed two-person economy with identical
concave consumption rule $c(a)$, transferring an infinitesimal dollar from
wealth $a_H$ to $a_L<a_H$ gives $c'(a_L)-c'(a_H)\geq0$. That proves a
fixed-price sign only. Changing income risk, access to credit, or curvature
changes this calculation, while the matrix $F_z^{-1}$ can amplify or
reverse aggregate feedback. No universal ordering follows from an inequality
index alone.

The formulas respect fixed initial households and zero initial net transfers.
They do not balance future government budgets automatically; those must be
inside $F$. The original identification problem remains: measured micro
responses do not by themselves identify the fiscal, price and transition
Jacobians. A concrete next task is joint experimental identification of those
Jacobians with sensitivity bounds on unobserved entries, then propagation
through an explicitly bounded inverse. \textbf{Final status: UNRESOLVED.}

\subsection*{Source}
B. Moll, \emph{Lecture 11: Good Research Topics and Open Questions},
Distributional Macroeconomics lecture slides, archived at a 2019 URL,
PDF page 9 (aggregate consumption and inequality).
\url{https://benjaminmoll.com/wp-content/uploads/2019/07/Lecture11_2149.pdf}.

\end{document}
